% =============================================================
% TABLAS COMPARATIVAS — ESTADO DEL ARTE  (versión 3 — verificada)
% Cambios respecto a v2:
%   (1) Se agrega Skin-R1 y sus baselines en HAM10000 (Tabla 2 y 3 del paper).
%   (2) Se corrige MM-Skin: BLEU-4 = 22.04 era para patología;
%       para dermoscopía: BLEU-4 = 2.84 (SkinVL-MM).
%   (3) Se agregan baselines de clasificación de MM-Skin paper en HAM10000.
%   (4) Se incluye DenseNet121 de Kreouzi (2024) como referencia binaria.
%   (5) Se incorporan resultados experimentales de este trabajo.
%   (6) Se separa tabla de generación en: métricas automáticas y NER clínico.
%
% REQUISITOS EN EL PREÁMBULO:
%   \usepackage{booktabs}
%   \usepackage{graphicx}
%   \usepackage{xcolor}
%   \usepackage{colortbl}
%   \usepackage{multirow}
%
% BIBLIOGRAFÍA:
%   \addbibresource{referencias_sota.bib}  % o incluir en tu .bib principal
% =============================================================

\definecolor{thiswork}{RGB}{230, 242, 255}
\definecolor{bestval}{RGB}{200, 230, 200}

% ──────────────────────────────────────────────────────────────
%  TABLA 1 — Clasificación de lesiones dermatológicas (HAM10000)
% ──────────────────────────────────────────────────────────────
\begin{table}[htbp]
\centering
\caption{%
  Comparativa del estado del arte en clasificación de lesiones dermatológicas
  sobre el dataset \textbf{HAM10000} (7 clases).
  $^\dagger$~indica arquitectura basada en VLM\@.
  \textit{FT}: ajuste fino supervisado;
  \textit{ZS}: evaluación \textit{zero-shot};
  \textit{RL}: aprendizaje por refuerzo.
  ``---''~indica métrica no reportada.
  Las filas sombreadas corresponden a los resultados de este trabajo.
}
\label{tab:sota_clasificacion_ham}
\resizebox{\textwidth}{!}{%
\begin{tabular}{@{}llccrrrrr@{}}
\toprule
\textbf{Referencia}
  & \textbf{Arquitectura}
  & \textbf{Modo}
  & \textbf{Cls.}
  & \textbf{Acc.\ (\%)}
  & \textbf{AUC-ROC}
  & \textbf{F1}
  & \textbf{Sens.\ macro}
  & \textbf{Espec.\ macro} \\
\midrule

% ── CNN / ViT ────────────────────────────────────────────────
\multicolumn{9}{@{}l}{\small\itshape Enfoques basados en CNN / ViT} \\[2pt]

\cite{SkinHealthMate2024}
  & EfficientNetB1$+$B5 (ens.)
  & FT & 7
  & \textbf{93.00} & --- & \textbf{0.93}$^{(w)}$ / 0.81$^{(m)}$ & --- & --- \\

\cite{SkinSavvy22025}
  & EfficientNetV2~$+$~Swin-Transf.
  & FT & 7
  & 77.70$^{(\dagger\dagger)}$ & --- & 0.793$^{(m)}$ & --- & --- \\

\addlinespace[5pt]
% ── VLM ──────────────────────────────────────────────────────
\multicolumn{9}{@{}l}{\small\itshape Enfoques basados en VLM$\,^\dagger$} \\[2pt]

\cite{SkinR12025}$^\dagger$
  & Skin-R1 (Qwen2.5-VL-7B~$+$~RL)
  & FT+RL & 7
  & \textbf{87.86} & --- & \textbf{0.879} & --- & --- \\

\addlinespace[3pt]
\multicolumn{9}{@{}l}{\small\itshape Baselines evaluados en HAM10000 por~\cite{SkinR12025}} \\[2pt]

\cite{SkinR12025}$^\dagger$~(baseline)
  & LLaVA-Med-7B
  & ZS & 7
  & 62.16 & --- & 0.577 & --- & --- \\

\cite{SkinR12025}$^\dagger$~(baseline)
  & GPT-4V
  & ZS & 7
  & 55.21 & --- & 0.561 & --- & --- \\

\cite{SkinR12025}$^\dagger$~(baseline)
  & LLaVA-v1.6-13B
  & ZS & 7
  & 61.25 & --- & 0.466 & --- & --- \\

\cite{SkinR12025}$^\dagger$~(baseline)
  & SkinVL-PubMM (LLaVA-Med$+$LoRA)
  & FT & 7
  & 60.54 & --- & 0.311 & --- & --- \\

\cite{SkinR12025}$^\dagger$~(baseline)
  & MedGemma-4B
  & ZS & 7
  & 58.59 & --- & 0.474 & --- & --- \\

\addlinespace[3pt]
\multicolumn{9}{@{}l}{\small\itshape Baselines evaluados en HAM10000 por~\cite{MMSkin2025}} \\[2pt]

\cite{MMSkin2025}$^\dagger$~(baseline)
  & LLaVA-Med-v1.5-7B
  & FT & 7
  & 62.59 & --- & --- & --- & --- \\

\cite{MMSkin2025}$^\dagger$
  & SkinVL-Pub (LLaVA-Med$+$LoRA)
  & FT & 7
  & 63.46 & --- & --- & --- & --- \\

\midrule
% ── Este trabajo ─────────────────────────────────────────────
\multicolumn{9}{@{}l}{\small\itshape Este trabajo --- clasificadores supervisados$^{(\nu)}$} \\[2pt]

\rowcolor{thiswork}
(a) & DenseNet121
  & FT & 7
  & 84.56 & 0.9637 & 0.8517$^{(w)}$ & 77.61\% & 96.74\% \\

\rowcolor{thiswork}
(b) & ViT-B/16
  & FT & 7
  & \textbf{85.76} & \textbf{0.9653} & \textbf{0.8621}$^{(w)}$ & 79.72\% & \textbf{97.16\%} \\

\rowcolor{thiswork}
(a') & ConvNeXt-Tiny
  & FT & 7
  & 82.70 & 0.9605 & 0.8606$^{(w)}$ & 84.89\% & 97.07\% \\

\rowcolor{thiswork}
(b') & Swin-Tiny
  & FT & 7
  & 86.83 & 0.9511 & \textbf{0.8817}$^{(w)}$ & \textbf{87.04\%} & \textbf{97.56\%} \\

\addlinespace[5pt]
\multicolumn{9}{@{}l}{\small\itshape Este trabajo --- VLMs (evaluación \textit{zero-shot}, sin fine-tuning de clasificación)} \\[2pt]

\rowcolor{thiswork}
(c)$^\dagger$ & LLaVA-Med 7B
  & ZS & 7
  & 11.44 & --- & 0.030 & 15.61\% & 85.80\% \\

\rowcolor{thiswork}
(d)$^\dagger$ & Mini-VLM (CLIP ViT-B/16~$+$~Phi-3-mini)
  & ZS & 7
  & 5.12 & --- & 0.006 & 14.39\% & 85.71\% \\

\addlinespace[5pt]
\multicolumn{9}{@{}l}{\small\itshape Este trabajo --- VLMs (\textit{likelihood scoring} y \textit{Chain-of-Thought})$^{(\xi)}$} \\[2pt]

\rowcolor{thiswork}
(c.lik.zs)$^\dagger$ & LLaVA-Med 7B
  & ZS lik & 7
  & 3.66 & --- & 0.028 & 15.72\% & 85.64\% \\

\rowcolor{thiswork}
(c.lik.ft)$^\dagger$ & LLaVA-Med 7B (Stage 2 v1, LoRA)
  & FT lik & 7
  & 13.44 & --- & 0.098 & 16.30\% & 85.88\% \\

\rowcolor{thiswork}
(c.lik.ft.v2)$^\dagger$ & LLaVA-Med 7B (Stage 2 v2, LoRA)$^{(\rho)}$
  & FT lik & 7
  & 17.90 & --- & 0.198 & 13.72\% & 85.25\% \\

\rowcolor{thiswork}
(c.cot.zs)$^\dagger$ & LLaVA-Med 7B
  & ZS CoT & 7
  & 10.23 & --- & 0.019 & 14.29\% & 85.71\% \\

\rowcolor{thiswork}
(c.cot.ft)$^\dagger$ & LLaVA-Med 7B (Stage 2 v1, LoRA)
  & FT CoT & 7
  & 8.12 & --- & 0.029 & 11.86\% & 85.58\% \\

\rowcolor{thiswork}
(c.cot.ft.v2)$^\dagger$ & LLaVA-Med 7B (Stage 2 v2, LoRA)$^{(\rho,\sigma)}$
  & FT CoT & 7
  & 16.79 & --- & 0.059 & 17.87\% & 86.05\% \\

\rowcolor{thiswork}
(d.lik.zs)$^\dagger$ & Mini-VLM
  & ZS lik & 7
  & 18.90 & --- & 0.268 & 18.08\% & 87.08\% \\

\rowcolor{thiswork}
(d.lik.ft)$^\dagger$ & Mini-VLM (Stage 2 v1, LoRA)
  & FT lik & 7
  & \textbf{66.60}$^{(\pi)}$ & --- & 0.536$^{(\pi)}$ & 14.21\% & 85.72\% \\

\rowcolor{thiswork}
(d.lik.ft.v2)$^\dagger$ & Mini-VLM (Stage 2 v2, LoRA)$^{(\rho,\pi)}$
  & FT lik & 7
  & 66.00 & --- & 0.532 & 14.09\% & 85.56\% \\

\rowcolor{thiswork}
(d.cot.zs)$^\dagger$ & Mini-VLM
  & ZS CoT & 7
  & 9.29 & --- & 0.022 & 12.37\% & 85.55\% \\

\rowcolor{thiswork}
(d.cot.ft)$^\dagger$ & Mini-VLM (Stage 2 v1, LoRA)
  & FT CoT & 7
  & 11.53 & --- & 0.040 & 14.13\% & 85.74\% \\

\rowcolor{thiswork}
(d.cot.ft.v2)$^\dagger$ & Mini-VLM (Stage 2 v2, LoRA)$^{(\rho)}$
  & FT CoT & 7
  & 11.12 & --- & 0.026 & 14.51\% & 85.72\% \\

\bottomrule
\end{tabular}%
}
\vspace{4pt}
\begin{minipage}{\textwidth}
  \footnotesize
  $^{(w)}$~F1 ponderado (\textit{weighted}); $^{(m)}$~F1 macro.
  $^{(\dagger\dagger)}$~Exactitud balanceada reportada sobre dataset HAM10000
  aumentado y rebalanceado (no sobre HAM10000 original);
  el F1-macro de 0.793 corresponde al Swin-Transformer~\cite{SkinSavvy22025}.\\
  Los baselines de~\cite{SkinR12025} provienen de su Tabla~3
  (evaluación \textit{in-distribution} sobre HAM10000, 7 clases).
  Los baselines de~\cite{MMSkin2025} provienen de su Tabla~4 de clasificación supervisada.\\
  $^{(\nu)}$~Los modelos (a), (b), (a') y (b') fueron seleccionados mediante barrido de hiperparámetros
  optimizando F1-weighted en el conjunto de validación.
  Los modelos (a) y (b) usan recipe \textit{legacy} con W\&B sweep ($25{+}$ corridas).
  Los modelos (a') y (b') usan el \textit{recipe} moderno (RandAugment, RandomErasing,
  \textit{two-stage fine-tuning}, TTA) con sweep grid de 15 trials por modelo
  (sweep v4: variación de \texttt{lr}$\times$\texttt{weight\_decay}$\times$\texttt{RandAug\_m}).
  Mejores trials: ConvNeXt-Tiny~t06 (\texttt{lr=7.5e-5, wd=0.10, m=6}),
  Swin-Tiny~t08 (\texttt{lr=1.5e-4, wd=0.05, m=6}).
  AUC-ROC calculado con normalización de probabilidades post-TTA.\\[4pt]
  $^{(\xi)}$~\textit{Likelihood scoring}: ranking de las 7 clases por la log-probabilidad
  del completion del LLM (length-normalized, ensemble de 4 prompts).
  \textit{Chain-of-Thought (CoT)}: el modelo describe primero la lesión y emite
  un diagnóstico final, que se parsea para mapear a una clase.
  \textit{FT}: checkpoint Stage 2 v1 con LoRA fine-tuning sobre SkinCaRe (~7k pares imagen-texto).\\[4pt]
  $^{(\pi)}$~Mini-VLM FT \textit{likelihood} alcanza \textit{accuracy} 66.6\% pero con
  \textit{sensitivity macro} 14.2\%, indicando colapso hacia la clase mayoritaria (NV, ${\sim}67\%$
  del test). Es un artefacto del prior lingüístico del Stage 2 entrenado sobre SkinCaRe
  donde NV domina; no representa capacidad discriminativa real.
  La iteración v2 (\texttt{epochs}=10) no rompe este colapso (sens.\ macro $\approx 14.1\%$).\\[4pt]
  $^{(\rho)}$~\textit{Stage 2 v2}: re-entrenamiento del Stage 2 con \texttt{epochs}=10 (vs 3 en v1)
  e hiperparámetros del mejor trial del sweep (t04: \texttt{lr=2.495e-5}, \texttt{lora\_r=64},
  \texttt{lora\_alpha=64}). Diseñado para mejorar convergencia. Resultado: F1 sube +0.100 absoluto
  en LLaVA-Med FT lik (+0.098$\to$0.198), pero \textbf{el class collapse persiste} en los modelos
  competitivos (\textit{sens.\ macro} sigue $\le 1/7$). El cuello de botella es \textit{imbalance + objective},
  no convergencia; ver detalles en \texttt{Codigo/docs/resul/vlm\_clasificacion\_resul\_v3\_20260517.md}.\\[4pt]
  $^{(\sigma)}$~LLaVA-Med FT CoT v2 produce salida no parseable (\textit{UNKNOWN}) en 81\% de los casos
  (1223/1503). El \textit{sensitivity macro} reportado se calcula sobre los 280 válidos y no es
  comparable con los demás combos.\\[4pt]
  \textbf{Nota sobre comparabilidad:}
  Todos los modelos en esta tabla fueron evaluados en HAM10000 con 7 clases.
  Sin embargo, los protocolos de evaluación difieren:
  \cite{SkinHealthMate2024} y~\cite{SkinSavvy22025} usan el split ISIC 2018;
  \cite{SkinR12025} usa evaluación basada en prompts con parseo de respuesta;
  este trabajo usa el split estándar 70/15/15 para clasificadores y prompt
  de clasificación directa para VLMs.
  La comparabilidad entre clasificadores supervisados y VLMs evaluados
  vía prompt es inherentemente limitada.
\end{minipage}
\end{table}


% ──────────────────────────────────────────────────────────────
%  TABLA 1b — Generalización OOD: evaluación en HIBA
% ──────────────────────────────────────────────────────────────
\begin{table}[htbp]
\centering
\caption{%
  Evaluación de generalización fuera de distribución (\textit{out-of-distribution}, OOD)
  sobre el dataset \textbf{HIBA} ($n=1616$), compuesto por imágenes dermoscópicas
  de pacientes argentinos con fototipos III--VI.
  Ningún modelo de esta tabla fue entrenado con datos HIBA.
  El mapping de clases HIBA$\to$HAM10000 es 1:1 para 6 de las 7 categorías;
  SCC se mapea a AKIEC (158 imágenes afectadas), lo que introduce ruido en esa clase.
  ``---''~indica que el modelo no fue evaluado sobre HIBA.
  Las filas sombreadas corresponden a resultados de este trabajo.
}
\label{tab:hiba_ood}
\resizebox{\textwidth}{!}{%
\begin{tabular}{@{}llcrrrrr@{}}
\toprule
\textbf{Modelo}
  & \textbf{Arquitectura}
  & \textbf{Recipe}
  & \textbf{Acc.\ (\%)}
  & \textbf{F1-w}
  & \textbf{AUC-ROC}
  & \textbf{Sens.\ macro}
  & \textbf{Espec.\ macro} \\
\midrule

\multicolumn{8}{@{}l}{\small\itshape Clasificadores supervisados (entrenados en HAM10000)} \\[2pt]

\rowcolor{thiswork}
(a) & DenseNet121   & legacy$^{(\nu)}$ & 56.56 & 0.5402 & 0.8372 & 43.96\% & 91.79\% \\

\rowcolor{thiswork}
(b) & ViT-B/16      & legacy$^{(\nu)}$ & 54.02 & 0.5241 & 0.8250 & 43.51\% & 91.56\% \\

\rowcolor{thiswork}
(a') & ConvNeXt-Tiny & sweep v4$^{(\nu)}$ & 47.15 & 0.5235 & 0.8127 & 49.56\% & 91.62\% \\

\rowcolor{thiswork}
(b') & Swin-Tiny     & sweep v4$^{(\nu)}$ & \textbf{58.48} & \textbf{0.5958} & 0.7912 & 53.20\% & \textbf{92.83\%} \\

\bottomrule
\end{tabular}%
}
\vspace{4pt}
\begin{minipage}{\textwidth}
  \footnotesize
  $^{(\nu)}$~Ver nota Tabla~\ref{tab:sota_clasificacion_ham}: todos los modelos seleccionados
  mediante barrido de hiperparámetros optimizando F1-weighted en validación HAM10000.\\[4pt]
  \textbf{Observaciones:}
  (1)~Todos los modelos pierden entre 27--40 puntos de exactitud al pasar de HAM10000 a HIBA,
  evidenciando un \textit{domain shift} significativo atribuible a diferencias de fototipo,
  condiciones de adquisición y distribución de clases.
  (2)~Swin-Tiny (sweep v4) obtiene el mejor F1-weighted OOD (0.596), superando a su versión v3
  (0.575) en $+0.020$ pt, aunque pierde AUC ($0.827{\rightarrow}0.791$, $-0.036$) y sensibilidad
  en melanoma ($0.545{\rightarrow}0.486$, $-0.059$). El sweep v4 mejora el calibrado del argmax
  pero no la separación probabilística OOD.
  (3)~ConvNeXt-Tiny (sweep v4) \textbf{degrada todas las métricas HIBA} respecto a v3
  ($\Delta$Acc${=}{-}0.084$, $\Delta$F1${=}{-}0.065$, $\Delta$Sens-MEL${=}{-}0.055$),
  sugiriendo que los HPs óptimos para HAM10000 producen un modelo sobre-ajustado a la
  distribución de origen y menos robusto al \textit{domain shift}.
  (4)~La disociación entre AUC y exactitud sigue presente, pero el sweep v4 reduce la AUC OOD
  en ambos modelos, indicando que la optimización \textit{in-distribution} no transfiere
  uniformemente a OOD.
  (5)~Ningún trabajo del SOTA reporta resultados en HIBA; esta tabla constituye
  la primera evaluación sistemática de modelos HAM10000 sobre población argentina.
\end{minipage}
\end{table}


% ──────────────────────────────────────────────────────────────
%  TABLA 2 — Clasificación con contexto adicional (otros datasets)
% ──────────────────────────────────────────────────────────────
\begin{table}[htbp]
\centering
\caption{%
  Referencia complementaria: resultados de clasificación con
  \textbf{modelos similares} a este trabajo pero en \textbf{otros datasets}.
  No directamente comparables con HAM10000 (7 clases), pero proporcionan
  contexto sobre las capacidades de las arquitecturas empleadas.
}
\label{tab:sota_clasificacion_otros}
\resizebox{\textwidth}{!}{%
\begin{tabular}{@{}lllccrrrr@{}}
\toprule
\textbf{Referencia}
  & \textbf{Arquitectura}
  & \textbf{Dataset}
  & \textbf{Modo}
  & \textbf{Cls.}
  & \textbf{Acc.\ (\%)}
  & \textbf{AUC-ROC}
  & \textbf{F1}
  & \textbf{Sens.} \\
\midrule

\cite{Kreouzi2024}
  & DenseNet121
  & DermNet (bin.)
  & FT & 2
  & 92.30 & 0.951 & --- & --- \\

\cite{EnsembleDL2025}
  & DenseNet121$+$ResNet50$+$EffNetB3
  & ISIC (bin.)
  & FT & 2
  & 97.67 & --- & 0.976$^{(\alpha)}$ & 0.963$^{(\alpha)}$ \\

\cite{PanDerm2024}
  & PanDerm (ViT-Large)
  & 2.1M imág.\ / 11 inst.
  & FT & bin.
  & --- & \textbf{0.964} & --- & 0.893 \\

\cite{MPXVCNN2023}
  & ResNet34
  & Multi-dataset (mpox)
  & FT & 2
  & --- & 0.966 & --- & 0.830$^{(\beta)}$ \\

\cite{MpoxVLM2024}$^\dagger$
  & CLIP~$+$~ViT~$+$~LLaMA-2-7B
  & Mpox dataset
  & FT & 2
  & 90.38 & 0.952 & 0.857 & 0.844 \\

\cite{DermLIP2025}$^\dagger$
  & DermLIP (CLIP ViT-B/16)
  & Derm1M
  & ZS & múlt.
  & 56.07 & 0.731 & --- & --- \\

\cite{DermatoLlama2025}$^\dagger$
  & DermatoLlama (Llama 3.2-11B)
  & DermaSynth (sint.)
  & FT & múlt.
  & 91.00 & --- & --- & --- \\

\bottomrule
\end{tabular}%
}
\vspace{4pt}
\begin{minipage}{\textwidth}
  \footnotesize
  $^{(\alpha)}$~F1 y Sensibilidad corresponden a la clase NV (benigna);
  métricas globales del ensemble no reportadas~\cite{EnsembleDL2025}.\\
  $^{(\beta)}$~Sensibilidad del cohorte de validación cruzada~\cite{MPXVCNN2023}.
  AUC proviene del cohorte de testeo externo.\\
  $^\dagger$~Modelo basado en VLM\@.
  Estas referencias son complementarias y no directamente comparables
  con la clasificación multiclase (7 clases) en HAM10000
  debido a diferencias en datasets, número de clases y protocolos.
\end{minipage}
\end{table}


% ──────────────────────────────────────────────────────────────
%  TABLA 3 — Generación de reportes: métricas automáticas
% ──────────────────────────────────────────────────────────────
\begin{table}[htbp]
\centering
\caption{%
  Comparativa del estado del arte en generación automática de reportes
  clínicos dermatológicos.
  Las métricas de texto están en la escala indicada entre paréntesis.
  ``---''~indica métrica no reportada.
  Las filas sombreadas serán completadas con los resultados experimentales
  (Nivel~1 pendiente de ejecución).
}
\label{tab:sota_reportes}
\resizebox{\textwidth}{!}{%
\begin{tabular}{@{}lllrrrrrl@{}}
\toprule
\textbf{Referencia}
  & \textbf{Arquitectura}
  & \textbf{Dataset (eval.)}
  & \textbf{BLEU-4}
  & \textbf{ROUGE-1}
  & \textbf{ROUGE-L}
  & \textbf{METEOR}
  & \textbf{BERTScore}
  & \textbf{Eval.\ clínica} \\
\midrule

% ── Evaluados sobre MM-Skin (dermoscopía) ───────────────────
\multicolumn{9}{@{}l}{\small\itshape Evaluados sobre MM-Skin --- imágenes dermoscópicas (VQA)} \\[2pt]

\cite{MMSkin2025}
  & SkinVL-MM (LLaVA-Med-v1.5$+$LoRA)
  & MM-Skin (dermoscopía)
  & 2.84
  & ---
  & 19.53
  & 22.41
  & 57.76$^{(\zeta)}$
  & --- \\

\cite{MMSkin2025}
  & SkinVL-PubMM (LLaVA-Med$+$LoRA$+$púb.)
  & MM-Skin (dermoscopía)
  & 2.60
  & ---
  & 24.53
  & 20.99
  & 60.94$^{(\zeta)}$
  & --- \\

\cite{MMSkin2025}~(baseline)
  & LLaVA-Med-v1.5-7B
  & MM-Skin (dermoscopía)
  & 2.58
  & ---
  & 14.81
  & 22.73
  & 63.31$^{(\zeta)}$
  & --- \\

\addlinespace[5pt]
% ── Evaluados sobre MM-Skin (patología) ─────────────────────
\multicolumn{9}{@{}l}{\small\itshape Evaluados sobre MM-Skin --- imágenes de patología (VQA)} \\[2pt]

\cite{MMSkin2025}
  & SkinVL-MM
  & MM-Skin (patología)
  & 22.04$^{(\eta)}$
  & ---
  & 44.09
  & 42.70
  & 70.59$^{(\zeta)}$
  & --- \\

\addlinespace[5pt]
% ── Otros datasets / tareas ────────────────────────────────
\multicolumn{9}{@{}l}{\small\itshape Evaluados sobre otros datasets (referencia metodológica)} \\[2pt]

\cite{AutomatedReport2025}
  & LLaVA-v1.5-13B (KD$^{(\theta)}$)
  & Dermoscopía priv.\ (305 img.)
  & 55.59$^{(\iota)}$
  & 54.38
  & 38.28
  & ---
  & 0.9025$^{(\kappa)}$
  & --- \\

\cite{DermatoLlama2025}
  & DermatoLlama (Llama 3.2-11B)
  & DERM12345 (Eval1)
  & 65.59$^{(\lambda)}$
  & 93.49$^{(\lambda)}$
  & 93.44$^{(\lambda)}$
  & ---
  & ---
  & Acc.\ MCQA: 90.93\% \\

\cite{SkinGPT4_2024}
  & ViT~$+$~Q-Former~$+$~LLaMA-2-13B
  & Privado (52k~img.)
  & ---
  & ---
  & ---
  & ---
  & ---
  & 80.63\%~correctos \\

\cite{PeFoMed2024}
  & ViT-EVA~$+$~LLaMA-2-7B~$+$~LoRA
  & IU-Xray$^{(\mu)}$
  & ---
  & ---
  & 28.60
  & 15.70
  & ---
  & WASM: 2.79/4 \\

\midrule
\multicolumn{9}{@{}l}{\small\itshape Este trabajo} \\[2pt]

\rowcolor{thiswork}
(c)$^\dagger$
  & Mini-VLM (CLIP ViT-B/16~$+$~Phi-3-mini)
  & MM-Skin (test split)
  & --- & --- & --- & --- & --- & --- \\

\rowcolor{thiswork}
(d)$^\dagger$
  & LLaVA-Med 7B~$+$~LoRA~$(r{=}64)$
  & MM-Skin (test split)
  & --- & --- & --- & --- & --- & --- \\

\bottomrule
\end{tabular}%
}
\vspace{4pt}
\begin{minipage}{\textwidth}
  \footnotesize
  $^{(\zeta)}$~BERTScore reportado en escala 0--100 en MM-Skin~\cite{MMSkin2025}.
  Las demás métricas (BLEU-4, ROUGE-L, METEOR) también en escala 0--100.\\
  $^{(\eta)}$~BLEU-4 sobre el subconjunto de \textbf{patología} de MM-Skin.
  Incluido para referencia; \textbf{no comparable con dermoscopía}.\\
  $^{(\theta)}$~Modelo estudiante entrenado mediante destilación de conocimiento
  con GPT-4V como maestro~\cite{AutomatedReport2025}.\\
  $^{(\iota)}$~Reportado como SacreBLEU (escala 0--100); no equivalente a
  BLEU-4 estándar. ROUGE-1 y ROUGE-L en escala 0--100.\\
  $^{(\lambda)}$~Valores originales en escala 0--1 multiplicados por 100.
  Los valores de ROUGE inusualmente altos pueden reflejar consistencia
  estructural en datos sintéticos~\cite{DermatoLlama2025}.\\
  $^{(\kappa)}$~BERTScore F1 en escala 0--1.\\
  $^{(\mu)}$~PeFoMed genera reportes radiológicos (IU-Xray), no dermatológicos;
  incluido como referencia de la arquitectura VLM$+$LoRA~\cite{PeFoMed2024}.\\[4pt]
  \textbf{Comparabilidad directa:} Las filas bajo ``MM-Skin --- dermoscopía''
  son las únicas directamente comparables con \textit{este trabajo},
  ya que comparten dataset y modalidad de imagen.
  Los demás papers usan datasets privados o tareas distintas.\\[4pt]
  \textbf{Nota sobre MM-Skin VQA:} Los valores de dermoscopía provienen
  de la Tabla~3 de~\cite{MMSkin2025} (columna ``Dermoscopy'').
  Se recomienda verificar los valores exactos contra el paper original,
  ya que la tabla reporta métricas por tres modalidades
  (patología, clínica, dermoscopía) con valores sustancialmente distintos.
\end{minipage}
\end{table}


% ──────────────────────────────────────────────────────────────
%  TABLA 4 — Evaluación NER Clínico (Nivel 2)
% ──────────────────────────────────────────────────────────────
\begin{table}[htbp]
\centering
\caption{%
  Evaluación de calidad clínica de los reportes generados
  mediante NER (\textit{Named Entity Recognition}) --- Nivel~2.
  \textit{Clinical Entity Recall}: proporción de categorías
  terminológicas dermatológicas mencionadas.
  \textit{Diagnostic Consistency}: porcentaje de reportes donde
  el diagnóstico correcto (o sinónimo) aparece en el texto.
  Estas métricas no requieren reporte de referencia.
}
\label{tab:nivel2_ner}
\resizebox{\textwidth}{!}{%
\begin{tabular}{@{}lrr@{}}
\toprule
\textbf{Modelo}
  & \textbf{Clinical Entity Recall (\%)}
  & \textbf{Diagnostic Consistency (\%)} \\
\midrule
\rowcolor{thiswork}
LLaVA-Med 7B~$+$~LoRA~$(r{=}64)$
  & 45.69
  & \textbf{71.47} \\
\rowcolor{thiswork}
Mini-VLM (CLIP ViT-B/16~$+$~Phi-3-mini)
  & \textbf{60.68}
  & 31.34 \\
\bottomrule
\end{tabular}
}
\vspace{4pt}
\begin{minipage}{0.95\textwidth}
  \footnotesize
  Evaluación basada en diccionario dermatológico personalizado
  (categorías ABCDE, tipos de lesión, características clínicas).
  Backend SciSpaCy no disponible en el entorno de evaluación;
  con NER biomédico activo los valores de \textit{entity recall}
  podrían ser mayores.
  Evaluado sobre 1.503 reportes generados a partir del test set HAM10000.\\[4pt]
  \textbf{Interpretación:} LLaVA-Med produce reportes más orientados
  al diagnóstico final (alta consistencia, 71.5\%);
  Mini-VLM genera descripciones morfológicas más ricas
  (mayor \textit{entity recall}, 60.7\%) pero con menor
  precisión diagnóstica (31.3\%).
\end{minipage}
\end{table}


% ──────────────────────────────────────────────────────────────
%  TABLA 5 — Resumen posicional de este trabajo vs. SOTA
% ──────────────────────────────────────────────────────────────
\begin{table}[htbp]
\centering
\caption{%
  Posicionamiento de los modelos de este trabajo respecto al estado del arte
  en HAM10000 (7 clases). Se incluyen solo referencias directamente
  comparables (mismo dataset, misma tarea).
}
\label{tab:posicionamiento}
\resizebox{\textwidth}{!}{%
\begin{tabular}{@{}llrrrl@{}}
\toprule
\textbf{Modelo}
  & \textbf{Tipo}
  & \textbf{Acc.\ (\%)}
  & \textbf{F1}
  & \textbf{$\Delta$ vs.\ SOTA$^{(\star)}$}
  & \textbf{Observaciones} \\
\midrule

\cite{SkinHealthMate2024} EfficientNet Ens.
  & CNN (FT) & \textbf{93.00} & \textbf{0.93}$^{(w)}$ & --- & SOTA absoluto CNN \\

\cite{SkinR12025} Skin-R1
  & VLM (FT$+$RL) & 87.86 & 0.879 & $-$5.14 & SOTA absoluto VLM \\

\rowcolor{thiswork}
(b') Swin-Tiny [\textit{este trabajo}]
  & ViT mod.\ (FT, sweep v4) & 86.83 & \textbf{0.882}$^{(w)}$ & $-$6.17 & \textbf{Mejor de este trabajo} \\

\rowcolor{thiswork}
(b) ViT-B/16 [\textit{este trabajo}]
  & ViT (FT) & 85.76 & 0.862$^{(w)}$ & $-$7.24 & --- \\

\rowcolor{thiswork}
(a') ConvNeXt-Tiny [\textit{este trabajo}]
  & CNN mod.\ (FT, sweep v4) & 82.70 & 0.861$^{(w)}$ & $-$10.30 & Mejor F1 entre CNNs \\

\rowcolor{thiswork}
(a) DenseNet121 [\textit{este trabajo}]
  & CNN (FT) & 84.56 & 0.852$^{(w)}$ & $-$8.44 & --- \\

\cite{SkinSavvy22025} Swin$+$EfficientNetV2
  & CNN (FT) & 77.70$^{(\dagger\dagger)}$ & 0.793$^{(m)}$ & $-$15.30 & HAM10000 aumentado \\

\addlinespace[4pt]
\cite{MMSkin2025} SkinVL-Pub
  & VLM (FT) & 63.46 & --- & $-$29.54 & LLaVA-Med$+$LoRA \\

\cite{SkinR12025} LLaVA-Med-7B
  & VLM (ZS) & 62.16 & 0.577 & $-$30.84 & Baseline de Skin-R1 \\

\cite{SkinR12025} LLaVA-v1.6-13B
  & VLM (ZS) & 61.25 & 0.466 & $-$31.75 & Modelo mayor, F1 menor \\

\cite{SkinR12025} GPT-4V
  & VLM (ZS) & 55.21 & 0.561 & $-$37.79 & --- \\

\addlinespace[4pt]
\rowcolor{thiswork}
(c) LLaVA-Med 7B [\textit{este trabajo}]
  & VLM (ZS) & 11.44 & 0.030 & $-$81.56 & Colapso de clase (método \textit{generate}) \\

\rowcolor{thiswork}
(d) Mini-VLM [\textit{este trabajo}]
  & VLM (ZS) & 5.12 & 0.006 & $-$87.88 & Colapso de clase (método \textit{generate}) \\

\bottomrule
\end{tabular}%
}
\vspace{4pt}
\begin{minipage}{\textwidth}
  \footnotesize
  $^{(\star)}$~$\Delta$ calculado respecto al mejor resultado global
  (EfficientNet Ensemble, 93.00\%).
  $^{(w)}$~F1 weighted; $^{(m)}$~F1 macro.
  $^{(\dagger\dagger)}$~Sobre HAM10000 aumentado (ver nota Tabla~\ref{tab:sota_clasificacion_ham}).\\[4pt]
  \textbf{Hallazgos clave:}
  (1)~Swin-Tiny sweep v4 (86.83\%, F1${=}0.882$) y ViT-B/16 (85.76\%, F1${=}0.862$) de este
  trabajo se posicionan como los mejores modelos individuales (no ensemble) del SOTA
  comparable, superando a todos los VLMs evaluados, incluyendo los fine-tuneados con LoRA.
  (2)~La brecha de 6--7 puntos con el SOTA CNN (93\%) se debe principalmente
  al uso de ensembles y pesos \textit{Noisy Student} en~\cite{SkinHealthMate2024}.
  (3)~La optimización de HPs (sweep v4) cierra la brecha entre arquitecturas modernas y
  legacy: Swin-Tiny v4 (0.882) supera a ViT-B/16 legacy (0.862), y ConvNeXt-Tiny v4 (0.861)
  iguala a ViT-B/16 legacy.
  (4)~Los VLMs en zero-shot (método \textit{generate}) colapsan hacia clases
  mayoritarias, resultando en accuracy por debajo del azar (14.3\% para 7 clases).
  El \textit{likelihood scoring} y CoT no aportan mejora significativa.
  Este es el punto de partida de la Fase~2 del trabajo.
  (5)~Skin-R1 (87.86\%) demuestra que con supervisión clínica estructurada
  y RL, los VLMs pueden superar a los clasificadores supervisados individuales,
  estableciendo el techo de la Fase~2.
\end{minipage}
\end{table}
